% =============================================================================
% conclusions_revised.tex
% -----------------------------------------------------------------------------
% Drop-in replacement for the manuscript's Discussion and Conclusion sections.
% This file is \input{} into learning_inventory_control_policies_es.tex and uses
% only the manuscript's preamble (natbib \citet/\citep, \section/\paragraph, the
% existing \label targets sec:discussion and sec:conclusion). No \documentclass.
%
% It replaces the current
%   \section{Discussion and managerial insights}            (label sec:discussion)
%   \section{Conclusion, limitations, and future work}      (label sec:conclusion)
% blocks. The two \section commands below carry the same labels, so all internal
% \ref{sec:discussion} / \ref{sec:conclusion} cross-references resolve unchanged.
%
% NEW \cite keys used here that are NOT yet in the manuscript's \thebibliography
% (the parent must add the matching \bibitem entries -- full details are listed at
% the foot of this file in a comment block):
%   alvo2023networks        Alvo, Russo, Kanoria, Lee (2023/2025) -- HDPO / reliable policy opt.
%   temizoz2025dcl          Temizoz, Imdahl, Dijkman, Lamghari-Idrissi, van Jaarsveld (2025) -- Deep Controlled Learning, EJOR
%   maggiar2025structure    Maggiar, Andaz, Bagaria, Eisenach, Foster, Gottesman, Perrault-Joncas (2025) -- Structure-Informed DRL
%   xie2024vc               Xie, Ma, Xin (2024) -- VC Theory for Inventory Policies
%   mania2018ars            Mania, Guy, Recht (2018) -- Simple/augmented random search
%   vanvuchelen2020jrp      Vanvuchelen, Gijsbrechts, Boute (2020) -- PPO for the joint replenishment problem
%   temizoz2024zsg          Temizoz, Imdahl, Dijkman, Lamghari-Idrissi, van Jaarsveld (2024) -- zero-shot generalization
% =============================================================================

\section{Discussion and managerial insights}
\label{sec:discussion}

Synthesizing across the ten problems, four themes emerge, and we close each by positioning it against
the deep reinforcement learning (DRL) and evolution-strategies literatures.

\paragraph{A single generic recipe is portable.}
The same CMA-ES loop, unchanged, learns a scalar lost-sales order, a setup-cost order/no-order
decision, a structured dual-sourcing order pair, a pair of echelon order-up-to levels, a per-node
backorder target vector, and a price-reactive purchase quantity. For a practitioner this means one
method to maintain rather than a bespoke heuristic per problem, and retraining when costs or demand
change is a matter of rerunning a short simulation rather than re-deriving structure. This portability
is the property the DRL roadmap of \citet{boute2022deep} asks of a learned controller, and the
cross-problem A3C study of \citet{gijsbrechts2022drl} first demonstrated for a neural learner.
Relative to that line, our contribution is not a new champion algorithm but a different point on the
cost--complexity frontier: a gradient-free optimizer over a compact policy reaches comparable territory
without value-function bootstrapping, reward shaping, or an exploration schedule, and so transfers
across structurally distinct problems with no per-problem hyperparameter tuning. This is the
inventory-control counterpart of the broader finding that black-box search over policy parameters is a
scalable alternative to gradient-based DRL \citep{Salimans2017EvolutionLearning}, and that even
linear policies found by simple random search can be competitive on hard control benchmarks
\citep{mania2018ars}. Within operations, evolutionary search has chiefly been used to tune the
parameters of a \emph{fixed} heuristic; learning the ordering policy itself with evolution strategies,
and doing so across many inventory families, is the gap we address.

\paragraph{Action geometry is where the leverage lives.}
Across every problem the backbone width mattered far less than the decoder. On vanilla lost sales a
small tree or linear ordinal head suffices; under a fixed cost the gated decoders that separate
\emph{whether} from \emph{how much} are what avoid the never-order collapse; in dual sourcing the
capped-dual-index coordinates are what let the policy reach the near-optimal region a raw-quantity
decoder cannot express; and in multi-echelon the direct-level design is the difference between a
$\approx 14\%$ improvement and a policy more than $200\%$ worse. The managerial reading is that
choosing the control representation --- ideally the coordinate system of the best known heuristic ---
is a more reliable lever than enlarging the network. This is, we believe, the most transferable
finding in the paper, and it is precisely the lever the most recent DRL literature has converged on
from the opposite direction: rather than freeing a large network to discover structure, that work
\emph{injects} inventory structure into the learner. \citet{maggiar2025structure} build
structure-informed policy networks that embed monotonicity and other analytically known properties of
optimal policies; \citet{temizoz2025dcl} report that a model-based controlled-learning scheme reaches
optimality gaps of at most $0.2\%$ on lost-sales and perishable problems, far inside earlier A3C gaps;
and recent graph- and heuristic-parameterized agents make the parameters of a classical control the
action space itself. Our policy-owned decoder is the same idea expressed as a design choice rather
than a regularizer: we put the search directly in the heuristic's coordinate system and let a compact
model modulate it. The convergence of independent lines on ``structure in the policy, not capacity in
the network'' is, in our reading, the clearest signal of where the field is heading.

\paragraph{Compact, interpretable, retrainable policies are a practical asset.}
The learned policies carry tens to a few hundred parameters --- orders of magnitude fewer than
published DRL networks, which the roadmap and subsequent studies note are large, tuning-heavy, and
hard to interpret \citep{boute2022deep,gijsbrechts2022drl,kaynov2024owmr}. They reduce at run time to
a small matrix multiply or a shallow tree traversal, so they can be deployed in standard enterprise or
spreadsheet tooling and audited by non-specialists; where a structured heuristic underlies the decoder
(dual sourcing, serial and divergent multi-echelon, general-network backorder), the learned outputs
read directly as order-up-to levels and caps. Interpretability is now an explicit objective in
inventory RL rather than an afterthought, and small or structurally constrained policies are
increasingly favored for exactly the audit-and-deploy reasons above. Compactness also carries a
statistical dividend that matters for practice: structured inventory policy classes provably generalize
from limited data, with generalization error largely independent of horizon for base-stock classes and
only logarithmic for $(s,S)$ classes \citep{xie2024vc}. A policy with hundreds of parameters living in
a heuristic's coordinate system is squarely the kind of low-complexity class for which such guarantees
apply, which is reassuring when the training signal is a finite simulation rather than an asymptotic
average.

\paragraph{Honest match/beat/bound accounting is itself a contribution.}
A recurring concern in this literature is that DRL inventory results are difficult to reproduce and
compare: cost conventions, demand processes, and baselines differ across papers, and there is no
ImageNet-style standard benchmark, so reported ``wins'' are not always like-for-like
\citep{boute2022deep,alvo2023networks}. We have tried to answer this concern at the level of
\emph{discipline} rather than scale. Every environment is anchored to a published number before it
carries any learned-policy claim (Appendix~\ref{app:validation}), and every result is reported under
one of three honest verdicts. We \emph{beat} a comparator only when it is a heuristic and the margin
clears the held-out standard error under common random numbers (general-network backorder by over
$20\%$; ameliorating inventory's order-up-to gate by a wide margin; one-warehouse multi-retailer on the
instances with exploitable dynamic structure; the production/assembly/distribution network's own best
pairwise base-stock). We \emph{match}, and say so explicitly, when the comparator is provably optimal
or a near-optimal proxy: the dual-sourcing soft tree sits inside the capped-dual-index policy's own
$\le 0.11\%$ optimality band --- comfortably inside the band the published A3C learner does not reach
--- and the serial soft tree recovers the Clark--Scarf optimum to $+0.011\%$; neither is framed as a
win. And we report a \emph{gap to a bound} without ever calling it ``beaten'': the ameliorating-inventory
gap to the perfect-information LP is structural (a single scalar purchase lever against the bound's full
three-part hindsight issuance) and is not comparable to the $\sim\!3.5\%$ deep-RL gap of
\citet{pahrGrunow2025}, who act in the full action space. Two comparisons we flag as indicative rather
than strictly like-for-like --- the negligible sub-band dual-sourcing margins, and the multi-echelon
improvement computed against a different baseline and cost convention than the published A3C figures.
This accounting is deliberately conservative, and it is where we believe a portable ES study can be
most useful to a field with a reproducibility problem: not by topping a leaderboard, but by stating
exactly what kind of claim each number supports.

\section{Conclusion, limitations, and future work}
\label{sec:conclusion}

We have shown that evolution strategies --- specifically CMA-ES over compact, policy-owned decoders ---
learn inventory control policies that are competitive across ten inventory problems: lost sales,
fixed-cost lost sales, dual sourcing, divergent and serial multi-echelon, general-network backorder,
perishable inventory, one-warehouse multi-retailer distribution, ameliorating inventory, and a
production/assembly/distribution network. Relative to DRL, CMA-ES needs few hyperparameters, supplies
its own exploration, and optimizes the realized cost directly, so the same loop transfers across
problems; the resulting policies are orders of magnitude smaller than published DRL networks while
remaining interpretable and verifiable. The consistent lever is action geometry: with the decoder
shaped like the relevant heuristic's coordinate system, a few-hundred-parameter policy matches proven
optima where they exist, beats strong heuristics where the comparator is a heuristic, and clears
published A3C on dual sourcing.

\paragraph{Limitations.}
Several caveats bound these claims and we state them plainly. (i)~Two comparisons are indicative rather
than strictly like-for-like: the dual-sourcing margins below the capped-dual-index proxy are negligible
and sit inside the heuristic's own optimality band, so the honest reading is a \emph{match}, not a win;
and the multi-echelon improvement is computed against a different baseline and cost convention than the
published A3C figures. (ii)~The high-penalty, autocorrelated MMPP lost-sales instances at the longest
lead times ($L=8,10$) remain won by classical heuristics, indicating a real gap for stationary compact
policies under regime-switching demand and deep pipelines. (iii)~We do not beat the strongest published
PPO on the one-warehouse multi-retailer instances, which remains the best learned policy on every row
there; our wins on that problem are over the tuned heuristic gate, not over PPO. (iv)~The
ameliorating-inventory gap is to a perfect-information \emph{upper bound} and is structural; it is
reported, never beaten. (v)~The production/assembly/distribution network result is a research result on
a faithful but \emph{not} literature-anchored environment: the learned policy beats the environment's
own best heuristic, not a published cost. (vi)~Methodologically, CMA-ES --- though far more
sample-efficient than A3C here --- is still data-hungry relative to settings with scarce historical
demand, and its population-based evaluation cost grows with the covariance dimension, so very large
parameterizations may need restricted, diagonal, or separable covariance structures.

\paragraph{Where the field is going, and how this fits.}
Three currents in recent inventory RL frame our future work. First, the field is moving toward
\emph{reproducible, optimum-anchored benchmarks}: open environments and comparisons to known or bounded
optima rather than to other learners \citep{alvo2023networks,temizoz2025dcl}. Our per-environment
literature validation and three-verdict accounting are a step in that direction, and contributing our
ten validated environments and policy forwards to such a shared suite is the natural next move.
Second, the most effective recent learners \emph{put inventory structure inside the policy} ---
structure-informed networks, base-stock-regularized objectives, and heuristic-parameterized action
spaces \citep{maggiar2025structure,temizoz2025dcl}. Our action-geometry result is the same principle
reached gradient-free; the obvious synthesis is to hand a sample-efficient gradient learner the same
policy-owned decoders, using CMA-ES (or warm-started from a heuristic, as we do for dual sourcing,
serial multi-echelon, and general-network backorder) to seed and the gradient learner to refine ---
which should directly attack the long-lead MMPP and the below-PPO gaps. Third, the frontier is turning
to \emph{generalist and limited-data} settings: generally-capable agents trained for zero-shot transfer
across instances, then paired with parameter estimation at deployment
\citep{temizoz2024zsg,vanvuchelen2020jrp}, and the favorable sample-complexity of structured policy
classes \citep{xie2024vc} suggests these compact decoders are exactly the right hypothesis class for
the data-scarce regime. Concrete next steps, in order of expected payoff, are: richer state-dependent
decoders for correlated and non-stationary demand; warm-starting from heuristics in every problem where
a structured control exists; scaling CMA-ES via separable or restricted covariance so larger networked
systems stay tractable; and amortizing one decoder across an instance family so that a single trained
policy generalizes with estimate-then-decide adaptation rather than per-instance retraining. The
through-line is unchanged from the body of the paper: in inventory control, the representation of the
\emph{action} is at least as consequential as the choice of \emph{learner}, and a compact, honest,
reproducible policy is often enough.

% =============================================================================
% NEW REFERENCES TO ADD (not already in the manuscript's \thebibliography).
% Add these \bibitem entries to the bibliography. None duplicate an existing key.
% Verified against arXiv / publisher pages, June 2026.
% -----------------------------------------------------------------------------
% \bibitem[Alvo et al.(2023)]{alvo2023networks}
% M.~Alvo, D.~Russo, Y.~Kanoria, and M.~Lee.
% Deep reinforcement learning for inventory networks: Toward reliable policy optimization.
% \emph{arXiv preprint arXiv:2306.11246}, 2023 (revised 2025).
%
% \bibitem[Temiz\"oz et al.(2025)]{temizoz2025dcl}
% T.~Temiz\"oz, C.~Imdahl, R.~Dijkman, D.~Lamghari-Idrissi, and W.~van Jaarsveld.
% Deep controlled learning for inventory control.
% \emph{European Journal of Operational Research}, 324(1):104--117, 2025.
%
% \bibitem[Maggiar et al.(2025)]{maggiar2025structure}
% A.~Maggiar, S.~Andaz, A.~Bagaria, C.~Eisenach, D.~Foster, O.~Gottesman, and
% D.~Perrault-Joncas.
% Structure-informed deep reinforcement learning for inventory management.
% \emph{arXiv preprint arXiv:2507.22040}, 2025.
%
% \bibitem[Xie et al.(2024)]{xie2024vc}
% Y.~Xie, W.~Ma, and L.~Xin.
% VC theory for inventory policies.
% \emph{arXiv preprint arXiv:2404.11509}, 2024.
%
% \bibitem[Mania et al.(2018)]{mania2018ars}
% H.~Mania, A.~Guy, and B.~Recht.
% Simple random search provides a competitive approach to reinforcement learning.
% \emph{arXiv preprint arXiv:1803.07055}, 2018.
%
% \bibitem[Vanvuchelen et al.(2020)]{vanvuchelen2020jrp}
% N.~Vanvuchelen, J.~Gijsbrechts, and R.~Boute.
% Use of proximal policy optimization for the joint replenishment problem.
% \emph{Computers in Industry}, 119:103239, 2020.
%
% \bibitem[Temiz\"oz et al.(2024)]{temizoz2024zsg}
% T.~Temiz\"oz, C.~Imdahl, R.~Dijkman, D.~Lamghari-Idrissi, and W.~van Jaarsveld.
% Zero-shot generalization in inventory management: Train, then estimate and decide.
% \emph{arXiv preprint arXiv:2411.00515}, 2024.
% =============================================================================
